\documentclass[11pt,a4paper]{article}
\usepackage[T1]{fontenc}
\usepackage[utf8]{inputenc}
\usepackage[english]{babel}
\usepackage{lmodern}
\usepackage[margin=27mm]{geometry}
\usepackage{microtype}
\usepackage[hidelinks]{hyperref}
\title{Telemetry Staleness and False Confidence in Predictive Humanoid Interaction Risk}
\author{Mateusz Pomianek}
\date{Version 0.1.0 -- 21 November 2026}
\begin{document}
\maketitle

\begin{abstract}
Stale or delayed sensor observations can make predictive risk estimates appear more certain than the underlying evidence warrants. How does controlled telemetry age and jitter change falsely confident low-risk decisions when a humanoid operates near people? The prospective design is as follows: Inject timestamped delay, jitter, missing packets and out-of-order ROS 2 observations into versioned offline HRI traces while preserving independent scene ground truth. The primary outcome is Rate of falsely confident low-risk decisions under predefined ground-truth hazards, stratified by telemetry age. The manuscript defines a research question and an evaluation strategy; it does not report experimental estimates.
\end{abstract}

\section{Introduction}
Stale or delayed sensor observations can make predictive risk estimates appear more certain than the underlying evidence warrants. The scientific contribution under consideration is a falsifiable comparison in which alternative explanations, measurement limitations and distinct hypotheses from related ARC Robo studies can be assessed independently. Both the source of observations and the assumptions needed to interpret results must be documented before any confirmatory comparison.

\section{Research question and hypothesis}
How does controlled telemetry age and jitter change falsely confident low-risk decisions when a humanoid operates near people? The working hypothesis is that the proposed approach improves the primary endpoint under an explicitly matched evidence and resource budget. This hypothesis is subject to prior-art review and prospective specification; non-superiority and negative findings must remain reportable.

\section{Materials and methods}
Inject timestamped delay, jitter, missing packets and out-of-order ROS 2 observations into versioned offline HRI traces while preserving independent scene ground truth. The experimental unit, independent ground truth, data inclusion criteria, version of software and sensors, and provenance of each dataset or scenario will be fixed in the protocol. Training, calibration and evaluation partitions will be isolated at the appropriate physical scene, sequence, participant or device level.

\subsection{Comparators and experimental controls}
Compare timestamp-agnostic inference, age-aware uncertainty inflation and conservative stale-data abstention using the same risk predictor. Model or policy hyperparameters must not be selected from final held-out results. The implementation will record relevant software versions, seeds, sensor configuration, timestamps and data rights.

\subsection{Primary endpoint}
Rate of falsely confident low-risk decisions under predefined ground-truth hazards, stratified by telemetry age. Denominators, operational thresholds, missing observations and failure cases will be defined before inferential analysis. Secondary metrics will not be substituted post hoc for the primary endpoint.

\section{Statistical and computational analysis plan}
Model risk-versus-data-age curves with scenario-disjoint intervals; report conditional risk coverage and tail-latency sensitivity. Any confirmatory statistical contrast will report uncertainty intervals, effective independent sample size and relevant failure strata. Where a physical validation or participant-based experiment is necessary, approval and feasibility assessment must precede data collection.

\section{Anticipated interpretation and limitations}
Timing perturbations in offline replay do not constitute evidence of a physical robot's fault response; the study must remain distinct from HRI-P04 calibration and HRI-P07 fallback assurance. Confounding from software versions, sensor conditions, scene selection and methodological overlap will be assessed explicitly. Claims about simulated behavior will not be generalized to physical safety or human behavior without matching evidence.

\section{Results}
Empirical results, tables and uncertainty intervals will be incorporated after execution of the documented evaluation protocol and independent verification of the input evidence.

\section{Discussion}
The planned comparison is intended to determine whether the proposed mechanism yields an identifiable gain over the stated controls. Interpretation must distinguish failure of a scientific hypothesis from an incomplete implementation, and should describe practical operating conditions under which conclusions remain valid.

\section{Conclusion}
This study defines a testable research problem, matched comparators and a primary evaluation target. A substantive empirical conclusion remains conditional on reproducible measurements and appropriate methodological review.

\section{Data and code availability}
The authoritative scientific source is maintained in the \texttt{MatPomGit/arc-hri-safety} repository. Dataset permissions, input manifests, source commits, evaluation scripts and permitted public artifacts will be identified in the final version. No participant records or proprietary instrumentation details are included here.
\end{document}
